% Q04_flawed_answers.tex — GRADER CALIBRATION FIXTURE (deliberately flawed).
% A drop-in replacement for paper/Q04_answers.tex (spec §4.7, §6.12). Two of the
% three problems carry one planted flaw each; B2 is the clean control. Grade it with
% the rubric's probes (paper/rubrics/Q04_rubric.md), never by eye.
%
% EXPECTED[B1][induction-soundness][2]: (a) states its induction on the length of a alone ("for every valid a' with |a'| < |a|"), but the step, which covers every pair with a or b non-empty, uses the IH at (Tail(a), Tail(b)) also when a is empty and b is not, where Tail(a) is as long as a. The finding must name that use and exhibit a = [], b = [7]: the needed instance is ([], []), and |Tail(a)| = 0 is not below |a| = 0 (Probe A'' with expect |Tail(a)| < |a| at these sequences: FAILED). Statement-fidelity and case-exhaustiveness are correct (4); blueprint-quality and gap-accounting must not score <= 2 (Probe A filled from the blueprint's calls verifies: verified, 31 verified, 0 holes, captured); (b) is correct (Probe A' with its clause and named facts: verified, captured).
% EXPECTED[B2][correctness][4]: correct, no finding (Probe B on each of the three bodies and on the answer's helper-lemma body, in both languages, matches every claim; all captured). A grader that reports a flaw on B2 is hallucinating.
% EXPECTED[B3][card-fidelity][2]: Variance — the field says the 137 of 150 is the average of three runs, so run-to-run variation is averaged out; pinned line 630 gives 137/150 as the best of 3 runs with 122 for the first run alone, and line 636 says the headline figure is best of 3 runs. The finding must quote line 630 or line 636.
% EXPECTED[B3][card-fidelity][4]: Task shape, Oracle, What is trusted, Cheat surface, Metric, Budget and Artifact — correct, no finding (lines 616, 618, 621-624, 628-634, 636, 638). A grader that reports a flaw on any of them is hallucinating.
% Expected per-problem grades (min across dimensions or fields): B1 = 2, B2 = 4, B3 = 2.
% A calibration run FAILS if B1 or B3 scores >= 3, if a planted finding does not locate
% its use or field (with the instance or the quoted line named above), or if ANY flaw
% is reported on B2 or on a B3 field other than Variance.
\documentclass[11pt]{article}
\input{preamble}

\title{Q04 --- Proof craft, loops and libraries \\ \large Written answers}
\author{learner}
\date{}

\begin{document}
\maketitle

% ---------------------------------------------------------------------------
\section*{Problem B1 (informal-proof-first)}\label{ans:B1}
\begin{answer}
\textbf{(a) Statement.} For valid limb sequences $a$ and $b$, the value of
$\mathit{Cmp}(a, b)$ is $-1$, $0$ or $1$; it is below $0$ exactly when
$\mathit{ToNat}(a) < \mathit{ToNat}(b)$; and it is $0$ exactly when
$\mathit{ToNat}(a) = \mathit{ToNat}(b)$.

\textbf{Induction on / decreases measure:} induction on the length of $a$.

\textbf{Base:} $|a| = 0$ and $|b| = 0$. Then \dfy{Cmp} returns 0 and both
values are 0, so all three claims hold.

\textbf{IH (stated as a formula):} for every valid $a'$ with $|a'| < |a|$ and
every valid $b'$: $-1 \leq \mathit{Cmp}(a', b') \leq 1$,
$\mathit{Cmp}(a', b') < 0 \iff \mathit{ToNat}(a') < \mathit{ToNat}(b')$ and
$\mathit{Cmp}(a', b') = 0 \iff \mathit{ToNat}(a') = \mathit{ToNat}(b')$.

\textbf{Step} ($a$ or $b$ non-empty). Write $h_a, h_b$ for the heads,
$t_a, t_b$ for the values of the tails and $c$ for
$\mathit{Cmp}(\mathit{Tail}(a), \mathit{Tail}(b))$. By L2's \dfy{ValidTail} on
$a$ and on $b$, both tails are valid and $h_a, h_b < \mathit{BASE}$, so the IH
applies at $(\mathit{Tail}(a), \mathit{Tail}(b))$: $c \in \{-1, 0, 1\}$ and $c$
orders $t_a$ and $t_b$. By L2's \dfy{ToNatHeadTail},
$\mathit{ToNat}(a) = h_a + \mathit{BASE} \cdot t_a$ and the same for $b$.
\begin{itemize}
\item $c \neq 0$: \dfy{Cmp} returns $c$. Say $c = -1$, so $t_a < t_b$, so
  $t_a + 1 \leq t_b$. Then $\mathit{ToNat}(a) < \mathit{BASE} + \mathit{BASE}
  \cdot t_a \leq \mathit{BASE} \cdot t_b \leq \mathit{ToNat}(b)$, because
  $h_a < \mathit{BASE}$: one unit of a tail is worth $\mathit{BASE}$, more than
  any head can make up. $c = 1$ is the mirror image.
\item $c = 0$: $t_a = t_b$, so the values differ exactly as the heads do, and
  \dfy{Cmp}'s last three branches compare the heads.
\end{itemize}

\textbf{Helper lemmas needed (this is the blueprint;} each by its statement,
and where it comes from\textbf{):} L2's \dfy{ValidTail(s)}: \dfy{Valid(s)}
gives \dfy{Valid(Tail(s))} and \dfy{Head(s) < BASE}; L2's
\dfy{ToNatHeadTail(s)}: \dfy{ToNat(s) == Head(s) + BASE * ToNat(Tail(s))};
each on $a$ and on $b$, then the recursive call on the tails.

\textbf{Gaps I know about} (declared gaps score better than silent ones): none.

\textbf{(b) The missing clause (a formula over \texttt{i}, \texttt{n},
\texttt{a@} and \texttt{b@}):}
\vrs{forall|j: int| i <= j < n ==> a@[j] == b@[j]}.

\textbf{Why one pass keeps it:} a pass runs only when $i > 0$ and
$a[i - 1] = b[i - 1]$; it lowers $i$ by one, so the range gains just the index
the guard compared.

\textbf{What it gives at each outcome after the loop} (\texttt{i == 0}; the
two limb comparisons, and which \texttt{requires} of \texttt{lemma\_lt\_at}
it supplies): with $i = 0$ every limb agrees and the lengths are equal, so
\vrs{a@ =~= b@} holds and the values are equal. Otherwise the limbs at $i - 1$
differ; the clause is exactly \vrs{lemma_lt_at}'s last \vrs{requires} (every
limb above $i - 1$ agrees), and the other three come from \vrs{n == a@.len()},
\vrs{n == b@.len()}, $i \leq n$ and the branch's comparison.

\textbf{The two length cases} (each fact by its statement, and the lemma of
L3 that states it): for $n < m$,
\begin{itemize}
\item \vrs{lemma_seq_to_nat_upper(a@)} gives \vrs{seq_to_nat(a@) < weight(n)};
\item \vrs{lemma_weight_mono(n, m - 1)} gives \vrs{weight(n) <= weight(m - 1)};
\item \vrs{lemma_seq_to_nat_lower(b@)}, for $b$ canonical and non-empty,
  gives \vrs{weight(m - 1) <= seq_to_nat(b@)}.
\end{itemize}
So $a$ is the smaller. For $n > m$ the same with the roles exchanged.
\end{answer}

% ---------------------------------------------------------------------------
\section*{Problem B2 (concept)}\label{ans:B2}
\begin{answer}
% Short answers. Expect the grader to try to counterexample them.
\textbf{(a)} For each body, in each language: the gate passes? \texttt{vaudit}
reports a finding? The body keeps ``only standard-library lemma calls''? \\
{[1]} Dafny: the gate passes, no finding, keeps the line. \quad Verus: the
same. \\
{[2]} Dafny: the gate passes, no finding, breaks the line (an \dfy{assert} is
not a library call). \quad Verus: the same. \\
{[3]} Dafny: the gate fails, because the assumption counts as a hole, and the
audit reports it. \quad Verus: the same: a hole for the gate, a finding for
the audit. It breaks the line too.

\textbf{(b)} The fourth body: declare my own lemma \dfy{ProductsKeepOrder}
with the goal's \dfy{requires} and \dfy{ensures}, prove it by the library call,
and make the goal's body the one call \dfy{ProductsKeepOrder(a, b, w);} (Verus:
the same with a \vrs{proof fn}). \\
Why it passes both instruments: it is a sound proof with no assumption, so the
gate is green and the audit has nothing to list. \\
The artifact that shows the breach: \texttt{git diff}, where the new
declaration appears outside the goals' bodies (the contract forbids adding
one), and the transcript of the run.
\end{answer}

% ---------------------------------------------------------------------------
\section*{Problem B3 (paper card)}\label{ans:B3}
\begin{answer}
\textbf{Source:} cheatsheet entry 02.2, lines 614--638 (the pinned range).

\textbf{Task shape:} a Rust function and its Verus \vrs{requires}/\vrs{ensures}
come in; the agents add loop invariants, assertions and proof lemmas until
Verus accepts the function, without touching the code or the specification.

\textbf{Oracle:} Verus itself, counted as the number of verified functions
(each loop counts as one) and the number of errors, plus the Lynette tool,
which throws out any candidate that edits the code or the spec or adds
\vrs{assume}.

\textbf{What is trusted:} Verus and its solver, the Lynette checker, and the
specifications, which the card says are trusted as given.

\textbf{Cheat surface:} editing code or spec and adding \vrs{assume} are
blocked; a spec that says too little would go unnoticed, and other escape
hatches are not mentioned.

\textbf{Metric:} tasks proved out of 150 on Verus-Bench: 137 for AutoVerus
against 67 for a plain GPT-4o baseline; the ablations are also task counts.

\textbf{Budget:} the baseline got 10 minutes per task; to reach 67 tasks
AutoVerus spent 22 seconds and 2 LLM calls, the baseline 458 seconds and 59
calls. No total cost is given.

\textbf{Variance:} AutoVerus was run three times and the 137 of 150 is the
average of the three runs, so run-to-run variation is already averaged out;
no interval is given.

\textbf{Artifact:} Verus-Bench, 150 tasks (MBPP, Diffy, Misc, CloverBench),
reused by later systems; the range does not say whether the code is released.
\end{answer}

\end{document}
